% EXCERPT: original lines 1-158
\section{Two-walk scales and preliminary comparisons}\label{sec:fluctuations}

We now measure fluctuations in the second coordinate at first hits of
successive height layers. Common true levels provide fresh pair
continuations and, for independent environments, an iid sequence of
symmetric increments. Their truncated variance will determine the
fluctuation scale. We first compare all first-hit positions with this
sequence, then transfer the comparison between two forms of conditioning.
The resulting short-block estimate and endpoint spread bounds will be
used in Sections~\ref{sec:gaussian} and~\ref{sec:clipping}.

On a successful prefix, set
\[
Y_j=(X_{T_j}-X_0)_2.
\]
Choose a deterministic integer median \(b_j\) of \(Y_j\) under \(P_0^+\),
with \(b_0=0\), and write
\[
W_H=\max_{0\le j\le H}|Y_j-b_j|.
\]
All these variables record displacement from the indicated starting site.

\subsection{Common true cuts}

Common spatial regeneration levels for two walks, reached at possibly
different path times, are central to the shared- and independent-environment
comparison of Rassoul-Agha and Sepp\"al\"ainen~\cite[Section~7]{RassoulAghaSeppalainen2009}.
Their invariance principle assumes regeneration-time moments. Here we
establish the common-cut identities directly from finite words, using the
spatial first moment proved in Section~\ref{sec:geometry}.

For two starting sites \(x,y\) on the same layer, we distinguish the
independent conditioned law \(P_x^+\otimes P_y^+\) from the annealed law
\(P_{x,y}^{++}\) of two walks in the same environment, conditioned on both
no-drop events \(D\). Let \(P_{x,y}^{\rm sh}\) denote the latter pair law
before conditioning, and put
\[
p_2(x,y)=E[q(x)q(y)].
\]
The positivity of \(q\) established in Section~\ref{sec:geometry} gives a
uniform lower bound
\[
\underline p_2:=\inf_{x_1=y_1}p_2(x,y)>0.
\]
Indeed, for every \(c>0\),
\[
P\bigl(\min(q(x),q(y))<c\bigr)\le 2P(q(0)<c),
\]
so a sufficiently small \(c\) makes \(q(x)q(y)\ge c^2\) on an event of
uniformly positive probability. Product-field mixing also gives
\[
p_2(x,y)\longrightarrow p^2
\qquad\text{as }\|x-y\|\longrightarrow\infty,\quad x_1=y_1.
\]
To verify this directly, approximate \(q(0)\) in \(L^1(P)\) by a bounded
function of finitely many rows and translate the approximation to \(x\)
and \(y\). The approximating functions are independent once their
translated row sets are disjoint.

Each marginal of the shared conditioned law is dominated uniformly by
the corresponding single conditioned law:
\[
P_{x,y}^{++}(X^{(1)}\in A)
 =\frac{E[P_{x,\omega}(A\cap D)q(y)]}{p_2(x,y)}
 \le \frac{p}{\underline p_2}P_x^+(A),
\]
and likewise for the second walk. With superscripts distinguishing the
walks, define
\[
\Delta_j=Y_j^{(2)}-Y_j^{(1)},\qquad
Z_j=y_2-x_2+\Delta_j.
\]
A \emph{common true level} is a level that is true for both walks at their
respective first hits. Heights are measured from the common starting
layer, and level zero is included under either conditioned law.

\begin{lemma}[Common words and their continuation laws]
\label{fluct-common-words}
Under either conditioned pair law, common true levels occur infinitely
often. At each such level the translated pair continuation is fresh:
under independent environments it has law \(P_0^+\otimes P_0^+\), while
under a shared environment, from terminal sites \(x',y'\), it has law
\(P_{x',y'}^{++}\) before translation. The latter transition law depends
on the terminal pair of sites.

These continuation rules remain valid at finite stopping indices of the
chain of common words, when the information retained consists of the
traversed words and their path prefixes.
\end{lemma}

\begin{proof}
Work first with the raw pair laws. Test a new integer level above both
explored path pasts, pausing each walk at its first arrival there. All
departure rows revealed by these prefixes lie strictly below the
candidate level. The conditional probability that both arrivals are
true is \(p^2\) in independent environments and \(p_2\) at the arrival
sites in a shared environment. Both probabilities have a uniform
positive lower bound.

If the candidate fails, resume the two walks, for example one step of
each per round, until one drops below the candidate. This failure is
detected in finite time. Then test a level strictly above both explored
pasts. Each walk tends to positive infinite height almost surely; this
also holds for the shared raw pair by its marginals. Thus, after every
detected failure, another candidate can be reached. The uniform success
bound proves the existence of a common true level above any prescribed
height. Conditioning on the two initial no-drop events preserves this
almost-sure conclusion.

For the continuation law, specify a pair of finite words from \(x,y\)
to their first hits \(x',y'\) of a common terminal level \(L>0\).
Call this word pair admissible when both prefixes avoid the initial
drop and every intermediate positive level has a witnessed failure of
common truth: at least one prefix has dropped below that level after its
first hit. These are exactly the required prefix conditions. Indeed,
once truth is imposed at \(L\), any failure of truth at an earlier level
must already have occurred in the specified prefixes.

Let \(A_{x,y}^{\rm sh}(w_1,w_2)\) be the raw annealed prefix weight of an
admissible word pair. Its quenched weight is the product of the two path
weights in the shared environment, including any repeated use of a row.
Every departure row lies below the terminal layer. The suffixes
constrained by truth use only rows at or above that layer, so site-row
independence gives the conditional word probability
\[
P_{x,y}^{++}\bigl((w_1,w_2)\text{ is the first common word}\bigr)
 =
 A_{x,y}^{\rm sh}(w_1,w_2)
 \frac{p_2(x',y')}{p_2(x,y)}.
\]
More generally, for any event \(B\) of the suffix pair, its joint
probability with this word pair is
\[
A_{x,y}^{\rm sh}(w_1,w_2)
\frac{p_2(x',y')}{p_2(x,y)}
P_{x',y'}^{++}(B).
\]
Here \(B\) is read with the new sites as starting sites. In particular,
the factor \(p_2(x',y')/p_2(x,y)\) is part of the shared word law; it
does not cancel in general.

In independent environments the corresponding terminal and initial
conditioning probabilities are both \(p^2\), so their ratio is one.
Translation of each walk separately then makes the relative suffix law
\(P_0^+\otimes P_0^+\), independent of the word just traversed. Enumerate
the countably many admissible finite word pairs and iterate these
factorizations to obtain the stated continuation rules. Summing over a
finite stopping index gives the same rules there. The information at
such an index does not include any unseen upper environment.
\end{proof}

In independent environments, let \(L\ge1\) be the integer width of a
common word and let \(S\) be its increment of \(\Delta\). The successive
pairs \((L_i,S_i)\), with their path decorations, are iid by
Lemma~\ref{fluct-common-words}. Exchanging the walks shows that \(S\) is
symmetric. In what follows, expectations of \(L,S\) refer to this
independent common-word law.


